\subsection{Hạn chót dưới trôi dạt: giải thích hậu nghiệm từ dữ liệu có sẵn}
\label{muc:haunghiem}

Định lý \ref{dl:hocong} dự đoán rằng dưới trôi dạt bất lợi, một hạn chót tính từ lời hứa đầu có
thể trở thành nguyên nhân làm rỗng tập phương án, còn trần hành vi thì không bao giờ. Khóa luận
\emph{chưa} chạy nhánh hạn chót nào. Nhưng từ các quỹ đạo đã chạy dưới trần hành vi có thể ước
lượng một hạn chót sẽ khả thi ở những ô nào: với mỗi ô, lấy cực đại theo khách của phần muộn cuối
$D^{+} = \max(0,\, p_T - p_0)$ trên chuỗi giờ trả đã công bố; ô được coi là khả thi dưới hạn chót
$\ell$ khi cực đại đó không vượt $\ell$.

\begin{table}[htbp]
\centering
\caption{Khả thi ước lượng của hạn chót một chiều từ quỹ đạo dưới trần hành vi 30 giây (giải
thích hậu nghiệm, không phải kiểm định)}
\label{bang:haunghiem}
\small
\begin{tabular}{@{}lp{4.2cm}rrrr@{}}
\toprule
\textbf{Trôi dạt} & \textbf{Cực đại $D^{+}$ theo ô (s)} & $\ell = 30$ & $\ell = 60$ & $\ell = 120$
& \textbf{Trần hành vi}\\
\midrule
$k = 0$ & $\le 28$ ở cả 20 ô & 20/20 & 20/20 & 20/20 & 0/20\\
$k = 1$ & 54, 60, 63, 78, 81, 85, 85, 87, 102 (9 ô) & 0/9 & 2/9 & 9/9 & 0/20\\
\bottomrule
\end{tabular}
\end{table}

Ba cột giữa là số ô ước lượng khả thi; cột cuối là số ô \emph{vô nghiệm do cam kết} đã quan sát
thật dưới trần hành vi (bảng \ref{bang:phanloai}). Bảng này phải đọc kèm ba giới hạn.
\emph{Thứ nhất}, chín ô ở $k = 1$ là các ô \emph{sống sót} dưới trần hành vi; mười một ô kia dừng
vì bất biến cửa sổ đón trước khi hết khung nên không có $D^{+}$ cuối --- mẫu bị chọn theo kết quả
của một nhánh can thiệp. \emph{Thứ hai}, quy tắc được dựng và ``kiểm'' trên cùng dữ liệu, nên đây là
\emph{giải thích hậu nghiệm}, không phải kiểm định. \emph{Thứ ba}, khung sáng chỉ có trôi bất lợi,
nên bảng chưa phân biệt được Định lý \ref{dl:hocong} với nhận định thô ``hạn chót hỏng khi có trôi
dạt''; Hệ quả \ref{hq:chieutroi} chỉ phân biệt được khi có cả chiều thuận lợi. Ngoài ra, vì
$D^{+} \le M^{+}$ nên ước lượng nghiêng về phía \emph{lạc quan} cho hạn chót: thực tế có thể còn vô
nghiệm nhiều hơn.

\subsection{Kiểm hồi tố cho cơ sở kết quả}
\label{muc:hoito}

Cùng lập luận cho một phép kiểm trên kết quả \emph{đã có}. Nếu Định lý \ref{dl:hocong} đúng thì
dưới cơ sở kết quả với trần 30 giây, một ô vô nghiệm khi và chỉ khi trên quỹ đạo gần dự phòng có
một khách mà tổng biến động nhìn thấy $V$ vượt 30 giây. Quỹ đạo của C1 trên panel tĩnh là gần dự
phòng, vì tổng thành phần hành vi của nó chỉ 128 giây trên $1.581$ khách, nên $V \approx X$.

Quy tắc đó đoán đúng tập ô vô nghiệm đã biết ở bảng \ref{bang:doicho}: các ô có khách với $V > 30$
giây dưới C1 là \emph{đúng} bốn ô 16-r2, 16-r4, 17-r1, 17-r3 --- trùng tuyệt đối. Các ô sát ngưỡng
cũng đúng: cực đại $V$ là $29{,}8$ giây (14-r3), $29{,}4$ giây (15-r2), $27{,}8$ giây (14-r2) và các
ô này khả thi; $32{,}0$, $32{,}5$, $34{,}2$, $38{,}7$ giây và các ô này vô nghiệm.

Phải gọi đúng tên kết quả này: đây là \emph{bằng chứng hồi tố}, vì tập ô vô nghiệm đã được biết
\emph{trước} khi viết quy tắc. Quy tắc không có tham số nào để chỉnh nên khó ``khớp quá mức'', nhưng
nó vẫn là kiểm tính nhất quán, không phải xác nhận độc lập. Nó cũng chỉ hợp lệ khi quỹ đạo tham
chiếu gần dự phòng: với trần 120 giây nhà điều hành được phép dời đáng kể nên $V$ không còn xấp xỉ
$X$, và quy tắc không được dùng để giải thích tập ô vô nghiệm ở trần đó.
